\section{ICLR 2027 Paper Checklist}
\label{app:paper-checklist}

(Reminder: this checklist is part of the paper and is required by the venue.)

\subsection*{Claims}
\begin{itemize}
  \item \textbf{Question 1}: For all models and algorithms presented, check if you
  include clear statements of the associated computational costs.
  [Yes] — the analytic traffic/compute model (Section~\ref{sec:analysis}) and the
  measured per-arm latencies carry explicit cost expressions and units.
  \item \textbf{Question 2}: For each experiment, check if you include clear
  statements of the goals, the hardware and software used, and the amount of
  compute used.
  [Yes] — every measured table states the accelerator (Ascend 910B2, CANN 8.5,
  torch\_npu 2.9), the launcher, process count, and the timing protocol.
  \item \textbf{Question 3}: For each experiment, check if you include the
  uncertainty/error bars on results and the number of trials.
  [Yes] — reported means $\pm$ standard deviation over three launches /
  five repeats where applicable, with per-launch raw JSON retained.
  \item \textbf{Question 4}: Check if you provide comparisons with multiple
  baselines.
  [Yes] — dense, static SVD, persistent low rank (multiple ranks), the
  near-equal-budget persistent control (persistent r401 at 48.2~MiB vs.\
  causal r384 at 48.1~MiB; the r512 cross-rank comparison is explicitly
  labeled not parameter-matched), and a shuffled-context control are
  all reported.
  \item \textbf{Question 5}: Check if you provide a breakdown of the methods
  by contributions (method, data, compute, infrastructure).
  [Yes] — the evidence ladder (Table~\ref{tab:evidencemap}) separates the
  measured mechanism result from analytic and open items.
  \item \textbf{Question 6}: Check if you discuss the limitations of your work.
  [Yes] — Discussion and the required-accelerator-validation list scope the
  causal-quality, branch-skipping, training-lifecycle, and end-to-end limits.
  \item \textbf{Question 7}: For any theoretical claim, check if you state both
  its assumptions and its complexity.
  [Yes] — Propositions are given with explicit algebraic conditions and
  complexity expressions.
\end{itemize}

\subsection*{Limitations}
\begin{itemize}
  \item \textbf{Question 8}: Check if you include clear statements of any
  limitations of your work.
  [Yes] — see Discussion; causal quality is measured (same-rank parity with a
  trained persistent control at r384, r401, and r512 from one coherent run
  family, strictly better than static SVD; at the near-equal-budget pair the
  generated arm carries a $+1.95\%$ relative-MSE penalty, so an adaptive rank
  allocation—not lifetime—remains the route to exact-budget equality), and the
  central limitation is that end-to-end benefit and branch-skipping dispatch
  remain open (mechanism verified on device; the measured v4 multi-core
  kernel resolves the beyond-L2 kernel-level deficit, up to $3.03\times$
  faster than v1 at $T{=}4096$).
\end{itemize}

\subsection*{Theory}
\begin{itemize}
  \item \textbf{Question 9}: Check if you include the complexity of any
  proposed algorithms.
  [Yes] — Section~\ref{sec:analysis} provides exact closed-form traffic,
  computation, and working-set expressions.
  \item \textbf{Question 10}: Check if you list the assumptions of any
  theoretical work clearly and describe the implications of relaxing them.
  [Yes] — the break-even analysis states its geometry assumptions and the
  proposition proofs note where the counts simplify.
\end{itemize}

\subsection*{Experimental Results}
\begin{itemize}
  \item \textbf{Question 11}: Check if you include the code, data, and
  instructions needed to reproduce the main experimental results.
  [Yes] — the harness, kernel source, contract-manifest generator, and
  per-run JSON (with SHA-256) are in the supplementary material and
  scheduled for public release at acceptance.
  \item \textbf{Question 12}: Check if you specify all the training details
  (hyperparameters, schedule, data splits, seeds).
  [Yes] — the oscillatory/launch configuration, ranks, seeds (2026--2029 for
  the quality families; earlier three-seed runs are labeled as such where
  retained), dtype, and frozen hashed data splits are declared in the
  protocol; every quality-table row is bound to its run family in
  \texttt{results/REGISTRY.md}.
  \item \textbf{Question 13}: Check if you report error bars for every
  experiment, and give a clear explanation of how they were computed.
  [Yes] — standard-deviation error bars with the number of launches stated
  per table.
  \item \textbf{Question 14}: Check if you provide the exact number of
  evaluation runs and hyperparameter searches performed.
  [Yes] — the number of launches, warmup/timed iterations, and repeats are
  stated in each table caption.
  \item \textbf{Question 15}: Check if you report the average of many
  runs, or the best run, with a clear explanation.
  [Yes] — means over repeated launches; the probe arms report medians over
  repeated full runs with the noise band stated.
  \item \textbf{Question 16}: When comparing methods, check if you discuss
  the choice of hyperparameters for all methods.
  [Yes] — ranks, controller geometry, dtype, and the matched parameter
  budget are specified for each arm.
\end{itemize}

\subsection*{Computational Experiments}
\begin{itemize}
  \item \textbf{Question 17}: State clearly whether the experiments were
  performed on a specific and corresponding hardware.
  [Yes] — all systems measurements are on Ascend 910B2 with the exact
  CANN/torch\_npu versions stated.
  \item \textbf{Question 18}: State the number of variants of the
  experimental system used in the experiments.
  [Yes] — the four-arm ladder plus the dense/static/persistent controls
  are enumerated.
  \item \textbf{Question 19}: State how much total compute time the
  experiments used.
  [Yes] — the protocol and evidence ladder report launch counts and the
  historical compute scope; per-run wall-clock is recorded in the retained
  logs.
  \item \textbf{Question 20}: Describe whether the computational experiments
  use text-based or synthetic data and, if constructed by the authors,
  whether they are provided.
  [Yes] — the quality protocol uses named corpora (GSM8K, MBPP, Alpaca,
  WikiText-103, HellaSwag) with hashed disjoint splits; the systems harness
  uses declared tensor geometry rather than external data.
\end{itemize}

\subsection*{Code}
\begin{itemize}
  \item \textbf{Question 21}: State clearly whether code is provided and,
  if so, where/how.
  [Yes] — code and per-run artifacts are in the supplementary material and
  scheduled for public release upon acceptance; see also the Data/Code
  availability statements.
  \item \textbf{Question 22}: State whether the code is runnable in a
  documented way (build, dependencies, commands).
  [Yes] — the kernel build (CMake), harness launchers, and reproduction
  commands are documented in the supplementary material.
  \item \textbf{Question 23}: State whether the code will be runnable on
  standard compute resources.
  [No / Not fully] — device measurements require an Ascend 910B2 + CANN 8.5
  host; the analytic and unit-test portions run on standard resources. We
  state this constraint explicitly.
\end{itemize}
